Pakistan's industrial-process activity data are distributed across manufacturing statistics,
trade records, company reports, product-use data and national greenhouse-gas reporting. We
develop a reproducible Industrial Processes and Product Use (IPPU) framework that links each
IPCC category to process-specific activity data and material balances. Categories lacking
sufficient evidence remain explicitly unestimated, with their data requirements documented.
The reference GHG module estimates cement calcination, ammonia production net of recovered carbon, iron and steel, and lubricant use. Its FY2023--24 total is \SI{25.535}{\mega\tonne} CO$_2$; categories lacking supported methods and activity remain unestimated. The accompanying air product has its own admitted process coverage, including ammonia, steel and soda ash; this does not turn every IPPU GHG omission into an estimate.

A source-period audit shows that \SI{48.011}{\mega\tonne} of cement previously attributed to
FY2023--24 corresponds to FY2021--22; PBS and State Bank statistics both report
\SI{39.566}{\mega\tonne} for FY2023--24. Accounting for \SI{3.792}{\mega\tonne} of
clinker exports and \SI{0.026}{\mega\tonne} of imports gives
\SI{33.440}{\mega\tonne} of inferred domestic clinker production. That Tier~1 inference
gives a cement process estimate of 17.389--21.504 Mt \COtwo{} across the IPCC
clinker fractions of 0.75 and 0.95, both ends priced as named national configurations.
The reference is neither end: it is the clinker the producers themselves report.

All sixteen producer filings independently report \SI{41.953}{\mega\tonne} of clinker,
125.5\% of the low-clinker Tier~1 inference and 101.4\% of the high case, which at the
common factor gives the reference \SI{21.816}{\mega\tonne} \COtwo{} (conditional 95\%
interval 20.227--23.452 Mt); the 0.75-fraction model is priced beside it as the national
configuration \emph{cement\_tier1\_model\_low}. The gap between the census and the
high-clinker inference indicates unresolved differences in reporting coverage, clinker
transfers and stock movements rather than a small residual activity gap. A matched calendar-2022
reconstruction gives 17.985--22.595 Mt \COtwo{}, compared with
\SI{27.906}{\mega\tonne} in Pakistan's official cement inventory. The full official IPPU
inventory is \SI{49.112}{\mega\tonne} \COe{}, including 19.400 Mt from chemical
industries.

These results show that clinker trade and reporting-period alignment materially affect cement
estimates, while company production data can potentially replace Tier~1 clinker assumptions.
The national annual core now covers July 2015--June 2025. This note retains the detailed FY2023--24 clinker and reporting comparisons. Plant-specific ammonia recovery, nitric-acid abatement and fluorinated-gas equipment data remain consequential limitations.
